\clearpage
\VerdierSemanticAnchor{front-editors-note}
\VerdierPhysicalPageAnchor{0008}
\chapter*{編者注}
\addcontentsline{toc}{chapter}{編者注}

\VerdierSemanticAnchor{front-editors-note-publication-history}
J.-L. Verdier を追悼して Université Paris VII が1989年10月19日に催した式典
において、M. Artin、P. Cartier、L. Illusie は J.-L. Verdier の博士論文を
出版する着想を得た。2年後、私も彼らに加わり、J. Le Potier と M. Broué の
支援を得て、本文を \TeX{} に入力するための資金を Paris VII（URA 212）と
ÉNS（DMI）から確保した。入力完了後、その結果を Verdier のタイプ原稿と入念
に照合した。しかし Verdier 自身が最終訂正を行っていなかったため、数学的な
校閲がなお必要であった。

\VerdierSemanticAnchor{front-editors-note-chapter-i-modifications}
この精密な校閲は、P. Cartier（第I章）、B. Keller（第III章）、そして全文に
ついては M. Zisman と私が、細心の注意を払って行った。列挙するほどでもない
多数の誤植を除けば、主な変更は次のとおりである。第I章の命題 (3.1.5) は同値
を主張していた。この逆向き（後では一度も使われない）が偽なので、同値を含意
に置き換えた。命題 (3.4.1) には符号の誤りがあり、第III章 ((1.3.4)、(b)、
(3.2.7.3)) の符号の誤りを引き起こしていた。そのため「反対三角圏」
（第II章、1.1.7）の定義を調整し、次の主張が真であり続けるようにした：
\emph{反対圏の導来圏は、元の導来圏の反対三角圏である}。

\VerdierSemanticAnchor{front-editors-note-chapter-ii-modifications}
第II章では、§1.2 の命題の順序を変更した。Verdier は執筆後、二つの完全三角
の直和が完全三角になるという主張は三角圏の他の公理から従うことに気づいて
いた。そのため公理の一覧からは削除したが、それに伴う再配列を行っていなかった
のである。さらに、後で使われず、検証にも失敗した段落 (2.3.6) と (2.3.7) を
削除した。定常スペクトル対象 (4.4.2) の定義はわずかに修正した。元の定義では
弱すぎて、後続の定理 (4.4.3) を証明できなかったためである。この変更が以後
の内容に影響しないことも確認した。

\clearpage
\VerdierPhysicalPageAnchor{0009}
\VerdierSemanticAnchor{front-editors-note-apparatus}
記号索引と用語索引は補完した。一方、参考文献は変更しなかった（ただし SGA
各巻については ``Lecture Notes'' での確定書誌情報を示した）。とりわけ、
参考文献番号13の欠落は補っていない。

\VerdierSemanticAnchor{front-editors-note-acknowledgements}
この企画の実現をさまざまな形で助けてくださった M. Broué、A. Bruguières、
P. Cartier、L. Gruson、L. Illusie、B. Keller、J. Le Potier、M. Zisman、
ならびに Université Paris VII、École Normale Supérieure、Institut des
Hautes Études Scientifiques に心から感謝したい。また、本書を辛抱強く
\textsc{AMS}-\TeX{} に入力してくださった C. Defosse と C. Gourgues にも
感謝する。

\VerdierSemanticAnchor{front-editors-note-memorial}
1994年12月に早逝した L. Doustaing にも、この感謝を伝えたかった。

\VerdierSemanticAnchor{front-editors-note-signature}
\begin{flushright}
パリ、1996年6月14日\\
Georges Maltsiniotis
\end{flushright}
